\begin{resumo}
O presente documento estabelece a modelagem estrutural e o projeto orientado a objetos do Sistema Inteligente de Gestão de Eventos Adversos baseado no método \textit{Global Trigger Tool} (SIGEA-GTT), desenvolvido na Universidade Federal de Sergipe (UFS). Fundamentado nos preceitos formais da especificação OMG UML 2.5.1, nos princípios de arquitetura de software SOLID e no catálogo de padrões de projeto corporativos e do GoF (\textit{Gang of Four}), o trabalho documenta exaustivamente todas as classes que estruturam o backend da aplicação em Java 21 e Spring Boot 3.3, bem como o modelo reativo do frontend em TypeScript e Angular 18+. O documento decompõe o sistema em camadas de software rigorosamente desacopladas: a Camada de Domínio, modelando as entidades JPA centrais (\texttt{User}, \texttt{GttModule}, \texttt{GttTrigger}, \texttt{HarmSeverity}, \texttt{AcademicClass}, \texttt{Activity} e \texttt{ActivitySubmission}) e seus respectivos objetos de valor embutidos e semiestruturados em JSONB; a Camada de Persistência, formalizando as interfaces de repositório baseadas em Spring Data JPA; a Camada de Serviços, encapsulando as regras de negócio clínicas, o controle transacional ACID e as validações institucionais de domínio exclusivo (\texttt{@academico.ufs.br}); a Camada de Apresentação REST, detalhando os controladores RESTful, o padrão \textit{Data Transfer Object} (DTO), mappers estruturais (MapStruct) e manipuladores globais de exceção; e a modelagem estrutural do frontend, ancorada em \textit{Standalone Components}, serviços singleton e gestão reativa de estado com \textit{Angular Signals}. São apresentados diagramas de classes em notação gráfica vetorial nativa TikZ para cada subsistema e camada arquitetural, acompanhados por dicionários descritivos de atributos, métodos, visibilidades e multiplicidades. Por fim, são consolidadas matrizes de rastreabilidade bidirecional mapeando as classes aos vinte e cinco requisitos funcionais do sistema, aos casos de uso, às tabelas físicas do banco de dados relacional e aos padrões de projeto aplicados, atestando a robustez estrutural, a alta coesão e o baixo acoplamento da solução.

\vspace{\onelineskip}
\noindent
\textbf{Palavras-chave}: Engenharia de Software. Diagrama de Classes. UML 2.5. Programação Orientada a Objetos. Spring Boot. Angular. SIGEA-GTT. DCOMP/UFS.
\end{resumo}
